\scanpage{133}
Notons que $\lambda'=\lambda^{1/i}$ \unsure{sera} aussi dans $\mathbb L$ (car $\lambda\overline\lambda=Q\deleted{\unsure{i}}$, et de $\unsure{\overline{\mathrm m}|=}$ pour ses conjugués\ldots), donc on voit aussi qu'il existe un $\lambda'$ \deleted{\unsure{tel que}} $\in\mathbb L$, avec $\deleted{\unsure{f}}\lambda'\overline{\lambda'}=|\lambda'|^2=Q$, tel que $Q(\lambda')$ soit \deleted{\unsure{p}} \inserted{(purement)} imaginaire et \uline{contenant} $K$; de plus, si $Q$ est entier, on peut prendre $\lambda'$ entier.
\mspage{132}
